\section*{Hoofdstuk \ref{ch:g5:mirrors-reflection} --- Spiegels en weerkaatsing}

\begin{solution}{exo:g5:mirrors-reflection:1}
\omterm{def:g1:light-and-shadows:source}{Licht} kaatst van een \omterm{def:g5:mirrors-reflection:reflection}{spiegel} terug en vertrekt precies even schuin
als het aankwam, aan de andere kant --- gelijke hoeken. Een straal die
recht van voren aankomt, kaatst recht terug langs zijn eigen weg.
\end{solution}

\begin{solution}{exo:g5:mirrors-reflection:2}
Het oppervlak van de muur is een landschap van piepkleine bobbels:
elke straal gehoorzaamt op zijn eigen hellinkje aan de kaatsregel, dus
verstrooien de stralen alle kanten op en lost het plaatje op. De
politoer van de \omterm{def:g5:mirrors-reflection:reflection}{spiegel} houdt alle kaatsen in de pas --- orde bewaard,
\omterm{prop:g5:mirrors-reflection:image}{beeld} bewaard.
\end{solution}

\begin{solution}{exo:g5:mirrors-reflection:3}
Achter het glas, \qty{2}{m} diep --- dus $2 + 2 = \qty{4}{m}$ van jou
vandaan.
\end{solution}

\begin{solution}{exo:g5:mirrors-reflection:4}
Het \omterm{prop:g5:mirrors-reflection:image}{beeld} loopt met dezelfde trage stap naar jou toe, en blijft
precies zo diep achter het glas als jij ervoor staat. De opening
sluit met twee keer jouw snelheid --- één stap van jou plus één van
het \omterm{prop:g5:mirrors-reflection:image}{beeld}.
\end{solution}

\begin{solution}{exo:g5:mirrors-reflection:5}
Een \omterm{def:g5:mirrors-reflection:reflection}{spiegel} verwisselt links en rechts, en een tweede
\omterm{def:g5:mirrors-reflection:reflection}{spiegel} verwisselt ze weer terug. Het woord staat vooraf omgekeerd,
zodat bestuurders vooruit het \emph{door hun
achteruitkijkspiegel} nog een keer verwisseld zien --- terug tot
gewoon AMBULANCE.
\end{solution}

\begin{solution}{exo:g5:mirrors-reflection:6}
Twee keer --- één keer bij elke gekantelde \omterm{def:g5:mirrors-reflection:reflection}{spiegel}. De bovenste
\omterm{def:g5:mirrors-reflection:reflection}{spiegel} vouwt het \omterm{def:g1:light-and-shadows:source}{licht} van het tafereel omlaag door de doos; de
onderste vouwt het naar buiten het oog in. Allebei staan ze onder een
halve rechte hoek, evenwijdig gekanteld.
\end{solution}

\begin{solution}{exo:g5:mirrors-reflection:7}
Kantel de \omterm{def:g5:mirrors-reflection:reflection}{spiegel} tot de teruggekaatste zonnestraal op het doel landt
--- de kaatsregel maakt de flits richtbaar; de Zon is de
\omterm{def:g5:everyday-energy:store}{energievoorraad} die het \omterm{def:g1:light-and-shadows:source}{licht} gratis levert. Een flits ter grootte
van een zakspiegel is kilometers ver te zien.
\end{solution}

\begin{solution}{exo:g5:mirrors-reflection:8}
Stil water is één glad oppervlak: de kaatsen blijven in de pas, en het
plaatje van de bergen overleeft --- een \omterm{def:g5:mirrors-reflection:reflection}{spiegel}. \omterm{def:g2:air-around-us:wind}{Wind} breekt het
oppervlak in ontelbare gekantelde vlakjes; elk vlakje kaatst volgens
de regel op zijn eigen schuinte, het leger raakt uit de pas, en het
\omterm{prop:g5:mirrors-reflection:image}{beeld} spat uiteen in dansend geglinster.
\end{solution}

\begin{solution}{exo:g5:mirrors-reflection:9}
Elke teruggekaatste straal bereikt het oog vanaf de \omterm{def:g5:mirrors-reflection:reflection}{spiegel}, maar het
oog verlengt hem recht achterwaarts, door het glas heen. Alle stralen
die van je neus zijn teruggekaatst, achterwaarts verlengd, kruisen
elkaar op één en dezelfde plek achter de muur --- dus bouwt het oog
daar een spookneus, en spookpunt voor spookpunt een hele spookjij. De
eensgezindheid van de stralen is wat de illusie scherp maakt.
\end{solution}

\begin{solution}{exo:g5:mirrors-reflection:10}
Monteer hem onder een halve rechte hoek met beide gangen (schuin over
de hoek, zoals de \omterm{def:g5:mirrors-reflection:reflection}{spiegels} van de periscoop). Een straal uit het
gangpad vouwt bij de \omterm{def:g5:mirrors-reflection:reflection}{spiegel} en loopt naar de deur --- en aangezien de
wegen van het \omterm{def:g1:light-and-shadows:source}{licht} beide kanten op werken, vouwt een straal vanaf de
deur langs precies dezelfde weg terug het gangpad in: elk uiteinde
ziet het andere in hetzelfde glas.
\end{solution}

\begin{solution}{exo:g5:mirrors-reflection:11}
Eén kaats verwisselt links en rechts; de periscoop kaatst het
\omterm{def:g1:light-and-shadows:source}{licht} twee keer, en de tweede \omterm{def:g5:mirrors-reflection:reflection}{spiegel} maakt de verwisseling van
de eerste ongedaan --- de omkering wordt teruggedraaid, en het
tafereel komt gewoon om en om aan.
\end{solution}
